\documentclass[11pt,a4paper]{article}
\usepackage[utf8]{inputenc}
\usepackage[vietnamese]{babel}
\usepackage[T5]{fontenc}
\usepackage{amsmath,amssymb,amsfonts}
\usepackage{geometry}
\usepackage{fancyhdr}
\usepackage{multicol}
\usepackage{xcolor}
\usepackage{tikz}
\usetikzlibrary{calc}

% Thiết lập khổ giấy A4 và căn lề tối ưu
\geometry{
	a4paper,
	left=15mm,
	right=15mm,
	top=15mm,
	bottom=18mm
}

% Cấu hình header & footer
\pagestyle{fancy}
\fancyhf{}
\renewcommand{\headrulewidth}{0pt}
\renewcommand{\footrulewidth}{0.4pt}
\rfoot{\footnotesize\color{subgray}Trang \thepage}
\lfoot{\footnotesize\color{subgray}Phiếu học tập Toán 10 -- Hệ thức lượng trong tam giác (GDPT 2018)}

% Định nghĩa bảng màu Infographic học thuật
\definecolor{primaryred}{HTML}{C62828}   % Đỏ son cho tiêu đề
\definecolor{inkblue}{HTML}{182A78}      % Xanh mực cho đường nét & khung
\definecolor{subgray}{HTML}{424242}      % Xám cho thông tin phụ
\definecolor{boxbg}{HTML}{F0F4FA}        % Nền xanh nhạt

% ------------------------------------------------------------------------------
% ĐỊNH NGHĨA CÁC MÔI TRƯỜNG & LỆNH TƯƠNG THÍCH CHUẨN GÓI ex_test & compile-dual
% ------------------------------------------------------------------------------
\newcounter{ex}
\newenvironment{ex}[1][]{%
	\refstepcounter{ex}%
	\par\smallskip\noindent\textbf{\color{inkblue}Câu \theex.} #1\ignorespaces
}{\par}

% \True là cờ đánh dấu phương án đúng (chuẩn ex_test)
\newcommand{\True}{}
\providecommand{\True}{}

% Lệnh \choice tự động điều chỉnh 1 dòng, 2 dòng hoặc 4 dòng tuỳ độ dài phương án
\newlength{\widthcha}
\newlength{\widthchb}
\newlength{\widthchc}
\newlength{\widthchd}

\newcommand{\choiceFour}[4]{%
	\par\smallskip
	\noindent\textbf{A.} #1\par
	\noindent\textbf{B.} #2\par
	\noindent\textbf{C.} #3\par
	\noindent\textbf{D.} #4\par\smallskip
}

\newcommand{\choiceTwo}[4]{%
	\ifdim\widthcha<0.46\linewidth
	\ifdim\widthchb<0.46\linewidth
	\ifdim\widthchc<0.46\linewidth
	\ifdim\widthchd<0.46\linewidth
	\par\smallskip\noindent
	\begin{minipage}[t]{0.48\linewidth}\textbf{A.} #1\end{minipage}\hfill
	\begin{minipage}[t]{0.48\linewidth}\textbf{B.} #2\end{minipage}\par\smallskip\noindent
	\begin{minipage}[t]{0.48\linewidth}\textbf{C.} #3\end{minipage}\hfill
	\begin{minipage}[t]{0.48\linewidth}\textbf{D.} #4\end{minipage}\par\smallskip
	\else\choiceFour{#1}{#2}{#3}{#4}\fi
	\else\choiceFour{#1}{#2}{#3}{#4}\fi
	\else\choiceFour{#1}{#2}{#3}{#4}\fi
	\else\choiceFour{#1}{#2}{#3}{#4}\fi
}

\newcommand{\choice}[4]{%
	\settowidth{\widthcha}{\textbf{A.} #1}%
	\settowidth{\widthchb}{\textbf{B.} #2}%
	\settowidth{\widthchc}{\textbf{C.} #3}%
	\settowidth{\widthchd}{\textbf{D.} #4}%
	\ifdim\widthcha<0.22\linewidth
	\ifdim\widthchb<0.22\linewidth
	\ifdim\widthchc<0.22\linewidth
	\ifdim\widthchd<0.22\linewidth
	\par\smallskip\noindent
	\begin{minipage}[t]{0.24\linewidth}\textbf{A.} #1\end{minipage}\hfill
	\begin{minipage}[t]{0.24\linewidth}\textbf{B.} #2\end{minipage}\hfill
	\begin{minipage}[t]{0.24\linewidth}\textbf{C.} #3\end{minipage}\hfill
	\begin{minipage}[t]{0.24\linewidth}\textbf{D.} #4\end{minipage}\par\smallskip
	\else\choiceTwo{#1}{#2}{#3}{#4}\fi
	\else\choiceTwo{#1}{#2}{#3}{#4}\fi
	\else\choiceTwo{#1}{#2}{#3}{#4}\fi
	\else\choiceTwo{#1}{#2}{#3}{#4}\fi
}

% Môi trường \choiceTF cho câu hỏi trắc nghiệm Đúng/Sai
\newcommand{\choiceTF}[4]{%
	\par\smallskip
	\noindent\begin{minipage}[t]{0.82\linewidth}\textbf{a)} #1\end{minipage}\hfill\fbox{\small\strut\ \textbf{Đ}\ \ \textbar\ \ \textbf{S}\ }\par\smallskip
	\noindent\begin{minipage}[t]{0.82\linewidth}\textbf{b)} #2\end{minipage}\hfill\fbox{\small\strut\ \textbf{Đ}\ \ \textbar\ \ \textbf{S}\ }\par\smallskip
	\noindent\begin{minipage}[t]{0.82\linewidth}\textbf{c)} #3\end{minipage}\hfill\fbox{\small\strut\ \textbf{Đ}\ \ \textbar\ \ \textbf{S}\ }\par\smallskip
	\noindent\begin{minipage}[t]{0.82\linewidth}\textbf{d)} #4\end{minipage}\hfill\fbox{\small\strut\ \textbf{Đ}\ \ \textbar\ \ \textbf{S}\ }\par\smallskip
}

% Lệnh điền đáp số cho phần trả lời ngắn
\newcommand{\shortans}[2][]{%
	\par\smallskip\noindent\textbf{Đáp số:} \dotfill\ \framebox[2.4cm]{\rule{0pt}{1.1em}#2}\par
}

% Định dạng tiêu đề các phân phần
\newcommand{\sectionheader}[2]{%
	\vspace{0.3cm}
	\noindent
	\begin{tikzpicture}
		\node[fill=primaryred!10, draw=primaryred, rounded corners=3pt, inner sep=4pt, text width=\linewidth-8pt] {
			\textbf{\color{primaryred}#1}\\[1pt]
			\footnotesize\textit{\color{subgray}#2}
		};
	\end{tikzpicture}
	\par\vspace{0.12cm}
}

\begin{document}
	
	% ------------------------------------------------------------------------------
	% TIÊU ĐỀ BÀI HỌC VÀ THÔNG TIN HỌC SINH
	% ------------------------------------------------------------------------------
	\begin{center}
		{\large \textbf{\color{primaryred}PHIẾU HỌC TẬP: HỆ THỨC LƯỢNG TRONG TAM GIÁC}}\\[1.5mm]
		{\small \textbf{\color{inkblue}Chương trình Toán 10 -- Khám phá kiến thức \& Luyện tập (GDPT 2018)}}\\[2mm]
		\noindent
		\begin{tikzpicture}
			\draw[color=inkblue!40, line width=0.8pt] (0,0) -- (\linewidth,0);
		\end{tikzpicture}\\[1mm]
		{\footnotesize \textbf{Họ và tên:} \dots\dots\dots\dots\dots\dots\dots\dots\dots\dots\dots\dots\dots\dots \quad 
			\textbf{Lớp:} \dots\dots\dots\dots\dots \quad 
			\textbf{Ngày:} \dots\dots/\dots\dots/20\dots\dots}
	\end{center}
	
	% KHUNG MỤC TIÊU HỌC TẬP
	\vspace{-0.1cm}
	\noindent
	\begin{tikzpicture}
		\node[fill=boxbg, draw=inkblue!30, rounded corners=4pt, inner sep=5pt, text width=\linewidth-10pt] {
			\textbf{\color{inkblue}Mục tiêu bài học:}\\[1pt]
			\footnotesize $\bullet$ Vận dụng thành thạo Định lý Côsin, Định lý Sin để tính các cạnh và góc trong tam giác.\\
			\footnotesize $\bullet$ Sử dụng linh hoạt các công thức tính diện tích tam giác $S$ và bán kính ngoại tiếp $R$, nội tiếp $r$.\\
			\footnotesize $\bullet$ Giải quyết bài toán mô hình hóa thực tế đo khoảng cách và chiều cao.
		};
	\end{tikzpicture}
	\par\vspace{0.15cm}
	
	% ==============================================================================
	% PHẦN I: TRẮC NGHIỆM NHIỀU LỰA CHỌN (10 CÂU)
	% ==============================================================================
	\sectionheader{PHẦN I. TRẮC NGHIỆM NHIỀU LỰA CHỌN (10 CÂU)}{Khoanh tròn vào chữ cái A, B, C hoặc D trước phương án trả lời đúng nhất.}
	
	\setcounter{ex}{0}
	\setlength{\columnsep}{18pt}
	\begin{multicols}{2}
		
		\begin{ex}
			Cho tam giác $ABC$ có $b = 8$, $c = 5$ và $\widehat{A} = 60^\circ$. Độ dài cạnh $a$ bằng bao nhiêu?
			\choice
			{\True $a = 7$}
			{$a = 49$}
			{$a = \sqrt{129}$}
			{$a = 39$}
		\end{ex}
		
		\begin{ex}
			Cho tam giác $ABC$ có độ dài ba cạnh $a = 6$, $b = 8$, $c = 10$. Số đo góc $\widehat{C}$ của tam giác bằng:
			\choice
			{$30^\circ$}
			{$45^\circ$}
			{$60^\circ$}
			{\True $90^\circ$}
		\end{ex}
		
		\begin{ex}
			Cho tam giác $ABC$ có $a = 10$ và $\widehat{A} = 30^\circ$. Bán kính $R$ của đường tròn ngoại tiếp tam giác $ABC$ bằng:
			\choice
			{$R = 5$}
			{\True $R = 10$}
			{$R = 20$}
			{$R = 10\sqrt{3}$}
		\end{ex}
		
		\begin{ex}
			Công thức nào sau đây tính đúng diện tích $S$ của tam giác $ABC$ theo hai cạnh $b, c$ và góc xen giữa $\widehat{A}$?
			\choice
			{$S = bc \sin A$}
			{$S = \frac{1}{2} bc \cos A$}
			{\True $S = \frac{1}{2} bc \sin A$}
			{$S = \frac{1}{2} bc \tan A$}
		\end{ex}
		
		\begin{ex}
			Cho tam giác $ABC$ có độ dài ba cạnh $a = 13$, $b = 14$, $c = 15$. Diện tích $S$ của tam giác $ABC$ bằng:
			\choice
			{\True $S = 84$}
			{$S = 42$}
			{$S = 168$}
			{$S = 21$}
		\end{ex}
		
		\begin{ex}
			Cho tam giác $ABC$ có độ dài ba cạnh $a = 13$, $b = 14$, $c = 15$. Bán kính đường tròn nội tiếp $r$ của tam giác $ABC$ bằng:
			\choice
			{$r = 8$}
			{$r = 6$}
			{\True $r = 4$}
			{$r = 2$}
		\end{ex}
		
		\begin{ex}
			Cho tam giác $ABC$ có $b = 7$, $c = 5$ và $\cos A = \frac{3}{5}$. Diện tích $S$ của tam giác $ABC$ bằng:
			\choice
			{$S = 21$}
			{\True $S = 14$}
			{$S = 28$}
			{$S = 10{,}5$}
		\end{ex}
		
		\begin{ex}
			Cho tam giác $ABC$ vuông tại $A$ có $a = 10$, $b = 6$, $c = 8$. Độ dài đường cao $h_a$ kẻ từ đỉnh $A$ bằng:
			\choice
			{\True $h_a = 4{,}8$}
			{$h_a = 2{,}4$}
			{$h_a = 5$}
			{$h_a = 9{,}6$}
		\end{ex}
		
		\begin{ex}
			Cho tam giác $ABC$ có $c = 6$, $\widehat{A} = 45^\circ$, $\widehat{B} = 75^\circ$. Bán kính đường tròn ngoại tiếp $R$ của tam giác $ABC$ bằng:
			\choice
			{\True $R = 2\sqrt{3}$}
			{$R = 4\sqrt{3}$}
			{$R = 6$}
			{$R = 3\sqrt{2}$}
		\end{ex}
		
		\begin{ex}
			Cho tam giác $ABC$ có $b = 10$, $c = 6$ và $\widehat{A} = 120^\circ$. Độ dài cạnh $a$ bằng:
			\choice
			{\True $a = 14$}
			{$a = 2\sqrt{19}$}
			{$a = 16$}
			{$a = \sqrt{76}$}
		\end{ex}
		
	\end{multicols}
	
	% ==============================================================================
	% PHẦN II: TRẮC NGHIỆM ĐÚNG - SAI (4 CÂU)
	% ==============================================================================
	\sectionheader{PHẦN II. TRẮC NGHIỆM ĐÚNG -- SAI (4 CÂU -- MỖI CÂU GỒM 4 Ý a, b, c, d)}{Học sinh xác định mỗi khẳng định là Đúng (Đ) hoặc Sai (S) và đánh dấu vào ô tương ứng.}
	
	\setcounter{ex}{0}
	
	\begin{ex}
		Cho tam giác $ABC$ có độ dài ba cạnh lần lượt là $a = 7$, $b = 8$, $c = 5$. Xét tính đúng/sai của các khẳng định sau:
		\choiceTF
		{\True Nửa chu vi của tam giác $ABC$ là $p = 10$.}
		{\True Diện tích của tam giác $ABC$ là $S = 10\sqrt{3}$.}
		{\True Bán kính đường tròn ngoại tiếp tam giác $ABC$ là $R = \dfrac{7\sqrt{3}}{3}$.}
		{Bán kính đường tròn nội tiếp tam giác $ABC$ là $r = 2\sqrt{3}$.}
	\end{ex}
	
	\begin{ex}
		Cho tam giác $ABC$ có hai cạnh $b = 6$, $c = 8$ và góc $\widehat{A} = 60^\circ$. Xét tính đúng/sai của các khẳng định sau:
		\choiceTF
		{\True Độ dài cạnh $a = 2\sqrt{13}$.}
		{\True Giá trị $\cos B = \dfrac{5\sqrt{13}}{26}$.}
		{\True Diện tích tam giác $ABC$ bằng $12\sqrt{3}$.}
		{Góc $\widehat{B}$ là góc tù.}
	\end{ex}
	
	\begin{ex}
		Cho tam giác $ABC$ có cạnh $c = 12$, $\widehat{A} = 45^\circ$ và $\widehat{C} = 30^\circ$. Xét tính đúng/sai của các khẳng định sau:
		\choiceTF
		{\True Số đo góc $\widehat{B} = 105^\circ$.}
		{\True Độ dài cạnh $a = 12\sqrt{2}$.}
		{\True Bán kính đường tròn ngoại tiếp $R = 12$.}
		{\True Diện tích tam giác $ABC$ bằng $36(\sqrt{3} + 1)$.}
	\end{ex}
	
	\begin{ex}
		Hai tàu cá cùng xuất phát từ bến $A$, đi thẳng theo hai hướng tạo với nhau một góc $60^\circ$. Tàu thứ nhất chạy với tốc độ $30\text{ km/h}$, tàu thứ hai chạy với tốc độ $40\text{ km/h}$. Xét tính đúng/sai của các khẳng định sau:
		\choiceTF
		{\True Quãng đường tàu thứ nhất đi được sau $2\text{ giờ}$ là $60\text{ km}$.}
		{\True Quãng đường tàu thứ hai đi được sau $2\text{ giờ}$ là $80\text{ km}$.}
		{\True Khoảng cách giữa hai tàu sau $2\text{ giờ}$ là $20\sqrt{13}\text{ km}$.}
		{Sau $2\text{ giờ}$, khoảng cách giữa hai tàu lớn hơn $80\text{ km}$.}
	\end{ex}
	
	% ==============================================================================
	% PHẦN III: CÂU HỎI TRẢ LỜI NGẮN (5 CÂU)
	% ==============================================================================
	\sectionheader{PHẦN III. CÂU HỎI TRẢ LỜI NGẮN (5 CÂU)}{Học sinh viết đáp số hoặc câu trả lời ngắn gọn vào ô trống tương ứng.}
	
	\setcounter{ex}{0}
	
	\begin{ex}
		Cho tam giác $ABC$ có $a = 13$, $b = 14$, $c = 15$. Tính độ dài đường cao $h_a$ kẻ từ đỉnh $A$.
		\shortans{12}
	\end{ex}
	
	\begin{ex}
		Cho tam giác $ABC$ có $a = 7$, $b = 5$, $c = 8$. Tính số đo góc $\widehat{A}$ (theo đơn vị độ).
		\shortans{$60^\circ$}
	\end{ex}
	
	\begin{ex}
		Cho tam giác $ABC$ có $b = 10$, $\widehat{B} = 30^\circ$. Tính độ dài bán kính đường tròn ngoại tiếp $R$ của tam giác $ABC$.
		\shortans{10}
	\end{ex}
	
	\begin{ex}
		Cho tam giác $ABC$ có $a = 4$, $b = 6$ và $\widehat{C} = 30^\circ$. Tính diện tích $S$ của tam giác $ABC$.
		\shortans{6}
	\end{ex}
	
	\begin{ex}
		Từ hai vị trí $A$ và $B$ trên mặt đất cách nhau $100\text{ m}$, người ta cùng nhìn thấy đỉnh $C$ của một tháp truyền hình dưới các góc nâng lần lượt là $30^\circ$ và $45^\circ$ ($A, B$ và chân tháp $H$ thẳng hàng trên mặt đất phẳng, $B$ nằm giữa $A$ và $H$). Tính chiều cao $CH$ của tháp truyền hình (làm tròn kết quả đến hàng phần mười theo đơn vị mét).
		\shortans{136,6 m}
	\end{ex}
	
	% ------------------------------------------------------------------------------
	% FOOTER - TỰ ĐÁNH GIÁ CỦA HỌC SINH
	% ------------------------------------------------------------------------------
	\vspace{0.3cm}
	\noindent
	\begin{tikzpicture}
		\node[fill=boxbg, draw=inkblue!20, rounded corners=4pt, inner sep=4pt, text width=\linewidth-8pt] {
			\footnotesize \textbf{Tự đánh giá mức độ hiểu bài:} \quad 
			$\square$ Cần cố gắng \quad $\square$ Đã hiểu \quad $\square$ Hoàn thành tốt \quad $\square$ Rất tự tin\\[2pt]
			\footnotesize \textbf{Điều em còn thắc mắc:} \dotfill
		};
	\end{tikzpicture}
	
\end{document}
